\documentclass[11pt]{article}

%=====================================================================
%  Attention as a smooth arrangement, and the predictive-coding
%  suboperad of Arr_Sm.
%
%  Stand-alone companion note to dynamic-algebra-potentials.tex.
%  Self-contained: it reproduces the paper's *rendered* notation
%  (\sarr, \lensob, \inpt/\outp, \cot, \Phiconf, \ltens, ...) with a
%  light article preamble, rather than importing the memoir machinery.
%=====================================================================

\usepackage[margin=1.05in]{geometry}
\usepackage{amsmath,amssymb,amsthm,mathtools}
\usepackage{xcolor}
\usepackage[colorlinks=true,linkcolor=blue!50!black,citecolor=orange!50!black,
            urlcolor=orange!50!black]{hyperref}
\usepackage[capitalize]{cleveref}        % must load after hyperref

\linespread{1.06}

%---------------- theorem environments ----------------
% Shared counter via aliascnt, so cleveref names each environment
% correctly (Remark/Observation/...), not all as "Proposition".
\usepackage{aliascnt}
\newcommand{\sharedthm}[3]{% #1=env, #2=Singular, #3=Plural
  \newaliascnt{#1}{proposition}%
  \newtheorem{#1}[#1]{#2}%
  \aliascntresetthe{#1}%
  \crefname{#1}{#2}{#3}}

\theoremstyle{plain}
\newtheorem{proposition}{Proposition}[section]
\crefname{proposition}{Proposition}{Propositions}
\sharedthm{observation}{Observation}{Observations}
\sharedthm{corollary}{Corollary}{Corollaries}

\theoremstyle{definition}
\sharedthm{definition}{Definition}{Definitions}
\sharedthm{construction}{Construction}{Constructions}
\sharedthm{example}{Example}{Examples}

\theoremstyle{remark}
\sharedthm{remark}{Remark}{Remarks}
\sharedthm{slogan}{Slogan}{Slogans}

%---------------- notation (mirrors the paper's output) ----------------
\newcommand{\rr}{\mathbb{R}}
\newcommand{\nn}{\mathbb{N}}
\newcommand{\cat}[1]{\mathcal{#1}}
\newcommand{\Cat}[1]{\mathbf{#1}}
\newcommand{\fun}[1]{\mathrm{#1}}
\newcommand{\Fun}[1]{\mathsf{#1}}        % named functor (sans-serif)

\newcommand{\Sm}{\mathrm{Sm}}
\newcommand{\arr}{\mathbb{A}\Cat{rr}}    % abstract arrangement operad
\newcommand{\sarr}{\arr_{\Sm}}           % smooth arrangement operad
\newcommand{\pc}{\mathbb{P}\Cat{C}}      % polynomial coalgebras (dynamics)

\newcommand{\mfd}{\Cat{Mfd}}
\newcommand{\rvect}{\Cat{RVect}}
\newcommand{\vect}{\Cat{Vect}}

\newcommand{\rv}[1]{\mathbf{#1}}         % reactive vector space \rv Q=(Q,\sharpR_Q)
\newcommand{\sharpR}{\sharp}

\newcommand{\inpt}[1]{#1^{-}}
\newcommand{\outp}[1]{#1^{+}}
\newcommand{\lensob}[1]{\binom{\inpt{#1}}{\outp{#1}}}
\newcommand{\ltens}{\boxtimes}           % lens tensor

\renewcommand{\cot}{\mathsf{cot}}        % cotangent functor Mfd -> Poly
\newcommand{\cotof}[1]{\cot(#1)}
\newcommand{\yon}{\mathcal{y}}
\newcommand{\potd}{\fun{d}}

\newcommand{\conf}{\mathrm{conf}}
\newcommand{\phase}{\mathrm{phase}}
\newcommand{\Phiconf}{\Phi_{\conf}}
\newcommand{\Phiphase}{\Phi_{\phase}}
\newcommand{\intg}{\mathfrak{i}}
\newcommand{\interpsm}{\mathrm{sm}}

\newcommand{\id}{\mathrm{id}}
\newcommand{\blank}{\mathord{\vcenter{\hbox{\scriptsize$\square$}}}}
\DeclareMathOperator{\Hom}{Hom}
\DeclareMathOperator{\softmax}{softmax}
\DeclareMathOperator{\sgn}{sgn}

\newcommand{\Wq}{W_{\mathrm q}}
\newcommand{\Wk}{W_{\mathrm k}}
\newcommand{\Wv}{W_{\mathrm v}}
\newcommand{\Wo}{W_{\mathrm o}}

%=====================================================================
\title{Attention as a smooth arrangement,\\
       and the predictive-coding suboperad of $\sarr$}
\author{David I.\ Spivak}
\date{June 24, 2026}

\begin{document}
\maketitle

\begin{abstract}
We place a self-attention head inside the smooth arrangement operad
$\sarr$ of the companion manuscript~\cite{DAP}: its query/key/value
weights are the reactive parameter space $\rv Q$, the softmax routing is
the backward (input) leg $\inpt f$ of the arrangement, an attention-pooled
read-out is the forward (output) leg $\outp f$, and a self-prediction loss
is the canonical potential $U$. Attention heads are then \emph{generators}
of a suboperad of $\sarr$, not a suboperad themselves; we describe the
suboperad they generate. Its interpretation under the configuration
integrator $\Phiconf$ is hierarchical predictive coding: the two legs of
the lens are exactly the top-down and bottom-up channels, the Moore
condition built into $\sarr$ is the predict-then-compare discipline, and
the \emph{type} of the potential, $U\colon Q\times\outp M\times\inpt N\to\rr$,
is already that of a prediction-error unit. The modern-Hopfield reading of
attention is the flat, one-level corner; the nested, bidirectional object
is the general one. We close with a cleaner generating set (a single
softmax-routing primitive over the linear suboperad) and a list of
open questions.
\end{abstract}

\tableofcontents

\section{The arrangement interface}\label{sec.interface}

We recall just enough of $\sarr$ to fix types. A \emph{box} is a lens
object $\lensob M=\binom{\inpt M}{\outp M}$, a pair of smooth manifolds:
$\inpt M$ is what the box receives, $\outp M$ is what it emits. A
\emph{smooth arrangement}
\[
  f\colon\lensob{M_1}\ltens\cdots\ltens\lensob{M_K}\longrightarrow\lensob N
\]
wires $K$ inner boxes into an outer box. Writing
$\outp M\coloneqq\prod_{i=1}^{K}\outp{M_i}$ and
$\inpt M\coloneqq\prod_{i=1}^{K}\inpt{M_i}$, the arrangement is a triple
\begin{equation}\label{eq.arr}
  f=\Bigl(\rv Q,\;\tbinom{\inpt f}{\outp f},\;U\Bigr),
\end{equation}
consisting of a reactive parameter space $\rv Q=(Q,\sharpR_Q)$ together
with three smooth maps
\begin{align}
  \outp f &\colon Q\times\outp M \longrightarrow \outp N, \label{eq.outf}\\
  \inpt f &\colon Q\times\outp M\times\inpt N \longrightarrow \inpt M, \label{eq.inf}\\
       U  &\colon Q\times\outp M\times\inpt N \longrightarrow \rr. \label{eq.U}
\end{align}
Here $Q$ is a finite-dimensional real vector space and
$\sharpR_Q\colon Q\to\vect(Q^{*},Q)$ is its reactive sharp, a smooth
choice of metric turning covectors into vectors; the potential $U$ is
the writer-monad decoration. Operadic substitution composes the lens
parts in the usual way and \emph{adds} potentials. The symmetric group
$\Sigma_K$ permutes the inner boxes.

Two features of~\eqref{eq.outf}--\eqref{eq.U} carry the whole story below
and are worth stating in isolation.

\begin{remark}[Directionality, and the Moore condition]\label{rmk.moore}
The output leg $\outp f$ does \emph{not} take an $\inpt N$ argument: the
outer emission is a function of the inner emissions $\outp M$ and the
parameters only. The input leg $\inpt f$ does take $\inpt N$: the outer
input is routed down to the inner inputs. So an arrangement reads
\emph{emissions} freely but cannot read its own \emph{incoming} message
instantaneously --- the Moore (as opposed to Mealy) condition. This is
exactly what makes feedback wirings well-posed once interpreted as
dynamics.
\end{remark}

\begin{remark}[The potential sees both ends]\label{rmk.bothends}
The potential's type~\eqref{eq.U}, $U\colon Q\times\outp M\times\inpt N\to\rr$,
gives it access to the inner emissions $\outp M$ and to the outer input
$\inpt N$ at once. We return to this in \cref{sec.semantics}; it is the
single most important type in the note.
\end{remark}

\section{Attention as a smooth arrangement}\label{sec.attention}

Fix an embedding dimension $E$, head dimensions $d_q=d_k=:d$ and
$d_v$, and let each inner box emit and receive vectors:
\[
  \outp{M_i}=\rr^{E},\qquad \inpt{M_i}=\rr^{E}\qquad(i=1,\dots,K).
\]
Both legs of every box are the embedding space $\rr^{E}$ (the residual
stream); the head space $\rr^{d_v}$ below is an internal bottleneck, never
an interface. We write a generic inner-emission tuple as
$h=(h_1,\dots,h_K)\in\outp M$, with $h_i\in\rr^{E}$.

\begin{construction}[An attention head]\label{con.head}
Take the reactive parameter space to have carrier
\[
  Q\;=\;\vect\!\bigl(\rr^{E},\,\rr^{d_q}\oplus\rr^{d_k}\oplus\rr^{d_v}\bigr)
        \;\oplus\;\vect\!\bigl(\rr^{d_v},\,\rr^{E}\bigr),
\]
i.e.\ a point of $Q$ is a quadruple $W=(\Wq,\Wk,\Wv,\Wo)$ of query, key,
value, and \emph{output} projections, and let $\sharpR_Q$ be a metric on
this weight space (\cref{rmk.optimizer}). Define the three maps
of~\eqref{eq.outf}--\eqref{eq.U} as follows.

\emph{Backward leg (the attention).} With the outer input
$n\coloneqq\inpt N$ adjoined to the pool as one extra key/value pair
$(\Wk n,\Wv n)$, set
\begin{equation}\label{eq.attn}
  \inpt f(W,h,n)_i
   \;=\;\Wo\!\!\!\sum_{j\in\{1,\dots,K\}\cup\{\star\}}\!\!\! A_{ij}\,\Wv h_j
   \;\in\;\rr^{E},
  \qquad
  A_{i\,\cdot}=\softmax\!\Bigl(\tfrac{1}{\sqrt d}\,\langle \Wq h_i,\,\Wk h_j\rangle\Bigr)_{j},
\end{equation}
where the index $\star$ contributes the key/value built from $n$. Thus
$\inpt f$ is precisely ``compute queries, keys, values; softmax the rows
of the Gram matrix; mix the values.''

\emph{Forward leg (the read-out).} Give the outer box a query $w_\star$ ---
either a learned vector carried in $\rv Q$ or one pooled from $h$ --- and
let it attend over the inner emissions:
\begin{equation}\label{eq.pool}
  \outp f(W,h)
   \;=\;\Wo\sum_{j=1}^{K}\softmax\!\Bigl(\tfrac{1}{\sqrt d}\,\langle w_\star,\,\Wk h_j\rangle\Bigr)_j\,\Wv h_j
   \;\in\;\outp N .
\end{equation}
By \cref{rmk.moore} the outer query may not depend on $n$.

\emph{Potential (the prediction).} Let
\begin{equation}\label{eq.Uattn}
  U(W,h,n)\;=\;\tfrac12\bigl\lVert\,h-\mathrm{pred}(W,n)\,\bigr\rVert^{2},
\end{equation}
the squared error between the bottom-up activity $h=\outp M$ and a
top-down prediction $\mathrm{pred}(W,n)$ assembled from the descending
context $n=\inpt N$.\qedhere
\end{construction}

So nothing in attention is bolted onto $\sarr$: the $\mathit{QKVO}$ weights
are the reactive space $\rv Q$, the softmax routing is the backward leg
$\inpt f$, the pooled read-out is the forward leg $\outp f$, and the
self-prediction loss is the potential $U$. The construction is a smooth
(indeed real-analytic) arrangement.

\begin{remark}[A notational hazard]\label{rmk.clash}
The companion paper writes $\rv Q=(Q,\sharpR_Q)$ for the reactive
\emph{parameter} space and $K$ for the \emph{arity} (number of inner
boxes). Attention's traditional letters $Q,K,V$ collide with both. We
keep the paper's $\rv Q$ and $K$ and use lowercase head dimensions
$d_q=d_k=d$, $d_v$ and roman projections $\Wq,\Wk,\Wv,\Wo$ throughout.
\end{remark}

\begin{remark}[The optimizer lives in the sharp]\label{rmk.optimizer}
The sharp $\sharpR_Q$ is a free choice of metric on weight space. A
constant Euclidean sharp gives plain gradient descent; a state-dependent
sharp gives a preconditioned / natural-gradient method. Thus the choice
of optimizer (e.g.\ a diagonal Adam-style preconditioner) is data of
$\rv Q$, not an addition to the dynamics.
\end{remark}

\begin{remark}[The coordinate-free data: QK and OV circuits]\label{rmk.circuits}
The auxiliary dimensions $d_q,d_v$ never reach the interface; they are rank
bounds. Eliminating them, a head is two low-rank maps on the embedding
space,
\[
\begin{aligned}
  \Theta&\coloneqq\Wq^{\top}\Wk\colon\rr^{E}\to\rr^{E}
   &&\text{(QK circuit; where to read)},\\
  \Omega&\coloneqq\Wo\Wv\colon\rr^{E}\to\rr^{E}
   &&\text{(OV circuit; what to write)},
\end{aligned}
\]
since $\langle\Wq h_i,\Wk h_j\rangle=h_i^{\top}\Theta h_j$ and the written
message is $\sum_j A_{ij}\,\Omega h_j$. A coordinate-free head is the pair
$(\Theta,\Omega)$ together with the softmax. In the predictive-coding
reading (\cref{sec.semantics}) $\Omega$ is the generative weight
(cause~$\mapsto$~predicted effect) and $\Theta$ its addressing. Tying
$\Wo=\Wv^{\top}$ and $\Wq=\Wk$ makes both symmetric --- the regime in which
the attention update is the exact gradient of an energy (\cref{rmk.corners}).
\end{remark}

\section{The outer box is a full participant}\label{sec.outer}

It is tempting to read $\outp f$ as a passive read-out. It is not: the
outer box participates in the attention on \emph{both} legs of the lens,
and the apparent inner/outer asymmetry is the operadic boundary, not a
privileging of the inner boxes.

\begin{itemize}
\item \emph{Up-leg} $\outp f\colon Q\times\outp M\to\outp N$ of
\eqref{eq.pool}: the outer box carries a \emph{query} and reads the inner
boxes' keys/values. This is attention pooling --- the outer box ``receives
through a learned query,'' exactly as an inner box does.
\item \emph{Down-leg} $\inpt f\colon Q\times\outp M\times\inpt N\to\inpt M$
of \eqref{eq.attn}: the outer input $\inpt N$ is projected to a
\emph{key/value} and dropped into the pool the inner boxes query over.
This is the route by which the outer box's content reaches the inner
boxes: cross-attention to descending context.
\end{itemize}

Thus inner boxes do self-attention to one another \emph{plus}
cross-attention to the outer input, while the outer box pools the inner
boxes. The only constraint is the one already noted.

\begin{observation}[No instantaneous self-read]\label{obs.feedthrough}
By \cref{rmk.moore} the up-leg $\outp f$ has no $\inpt N$ argument, so the
outer box's read cannot see its own descending message inside a single
arrangement. Any dependence of the outer read-out on the outer input is
mediated by a box's state across one interpretation step, i.e.\ across
depth/time. This is the well-posedness condition, and it is the
residual-stream-as-state mechanism.
\end{observation}

\begin{remark}[The split dissolves under nesting]\label{rmk.supertoken}
The outer box's query rides the up-wire ($\outp N$) while its key/value
ride the down-wire ($\inpt N$): what a box presents \emph{up} to its
parent (a pooled summary) differs from what it broadcasts \emph{down} to
its children (the context it received). In one arrangement this looks
asymmetric, but it dissolves under operadic substitution: when this box
becomes an inner box one level up, its up-read-out $\outp N$ \emph{is} the
single emission from which the parent computes its own $\mathit{QKV}$.
Every box carries unified $\mathit{QKV}$ at its own level. Nesting the
construction therefore yields \emph{hierarchical} attention, each
arrangement a ``super-token'' that pools its children into a summary
presented upward and broadcasts the context it receives downward
(cf.\ inducing-point set attention~\cite{SetTransformer} and the Perceiver
latent~\cite{Perceiver}).
\end{remark}

\section{The generated suboperad}\label{sec.suboperad}

\begin{observation}[Heads generate, they do not close]\label{obs.generate}
Composing two attention heads yields a two-layer attention network: still
a smooth arrangement (lens parts compose, potentials add), but not a
single head. Hence the attention heads are \emph{generators}; the
suboperad they generate is the class of finite attention networks.
Operadic composition is depth and the lens tensor $\ltens$ is width
(multi-head / parallel heads).
\end{observation}

Two standard architectural facts are forced by the operad structure
rather than assumed.

\begin{observation}[Composability supplies the output projection]\label{obs.residual}
To substitute one arrangement into a box of another, a received message
must be re-emittable at the next level: the message delivered to box $i$
must land in its emission space $\outp{M_i}=\rr^{E}$. The value lives in a
head space $\rr^{d_v}$, so composability requires a re-embedding
$\rr^{d_v}\to\rr^{E}$ --- which is exactly the output projection $\Wo$ of
\cref{con.head}. The residual stream therefore keeps constant width $E$
while $d_v$ remains a free internal rank; the operad demands that $\Wo$
\emph{exist}, not that $d_v=E$.
\end{observation}

\begin{observation}[Equivariance forces positional encoding]\label{obs.equivariance}
Operad morphisms respect the $\Sigma_K$-action permuting inner boxes, and
the softmax routing~\eqref{eq.attn} is $\Sigma_K$-equivariant. Order
information is thus invisible to a pure arrangement; recovering it requires
breaking the symmetry by hand, i.e.\ positional encoding.
\end{observation}

\begin{example}[The $0$-ary box learns the prior]\label{ex.zeroary}
At arity $K=0$ the attention sum in~\eqref{eq.attn}--\eqref{eq.pool} is
empty: attention contributes nothing. The output leg
$\outp f\colon Q\times 1\to\outp N$ is a parameterized constant, and the
potential~\eqref{eq.Uattn} drives that constant toward the data marginal.
A $0$-ary box still predicts --- but its predictive power is the
potential-driven learned prior, not attention.
\end{example}

\section{What the suboperad computes}\label{sec.semantics}

Interpreting a decorated arrangement is the functor
$\Phiconf\colon\sarr\to\pc$ of the companion paper (and its phase variant
$\Phiphase$): the potential is lifted to a force by the cotangent functor
$\cot$ and the sharp, and the resulting vector field is a dynamical
system. We read off what attention-with-prediction computes.

The two legs of \cref{sec.outer} are precisely the two channels of
predictive coding: the down-leg $\inpt f$ broadcasts a top-down prediction
to the children; the up-leg $\outp f$ returns a bottom-up summary; and the
Moore condition (\cref{obs.feedthrough}) is the predict-then-compare
discipline. The clincher is \cref{rmk.bothends}.

\begin{observation}[The potential is an error unit]\label{obs.erroreunit}
The type $U\colon Q\times\outp M\times\inpt N\to\rr$ is exactly that of a
predictive-coding error unit: $U$ compares the bottom-up activity $\outp M$
against the top-down context $\inpt N$, as in~\eqref{eq.Uattn}. The
framework supplies the error unit for free; ``predict the input'' is not an
add-on but the canonical decoration of this interface.
\end{observation}

\begin{slogan}\label{slo.pc}
Arrangements alone (no $U$, no $\Phi$) are a feedforward function class:
hierarchical attention, each arrangement a super-token
(\cref{rmk.supertoken}). Adding the potential $U$ and interpreting under
$\Phiconf$ makes the dynamics \emph{hierarchical predictive coding}:
predictions flow down the lens, errors flow up, and the state relaxes to
minimize $\textstyle\sum U$.
\end{slogan}

\begin{remark}[The integrator chooses the order]\label{rmk.order}
Predicting the input alone is served by $\Phiconf$, giving first-order,
relaxational predictive coding. Predicting the input \emph{and its
state-update} targets the momentum coordinate, which $\Phiconf$ discards;
the phase integrator $\Phiphase$ retains it, giving second-order,
oscillatory predictive coding. The choice of what to predict selects the
integrator.
\end{remark}

\begin{remark}[The flat corners]\label{rmk.corners}
The single-level instance recovers known objects. A flat pool with the
softmax routing~\eqref{eq.attn} and a quadratic energy is a modern
Hopfield network, whose energy-descent fixed-point update is
attention~\cite{HopfieldAllYou}. That identification of attention with
$-\nabla(\text{energy})$ is exact only for the symmetric, tied head
($\Wo=\Wv^{\top}$, $\Wq=\Wk$; \cref{rmk.circuits}); the untied transformer
is the non-gradient generalization, in which attention is the wiring
$\inpt f$ and the force a separate $-\nabla U$. A flat pool with a linear generative
model is classical Rao--Ballard predictive coding~\cite{RaoBallard}. The
nested, attention-routed object of \cref{slo.pc} is the general case;
these are its corners.
\end{remark}

\begin{remark}[Toward active inference]\label{rmk.active}
If the down-leg $\inpt f$ is closed back onto the environment --- so the
broadcast can \emph{act}, not merely predict --- the same potential-descent
dynamics become active inference / free-energy minimization in the sense
of~\cite{Friston}. The construction here is the perceptual-inference half.
\end{remark}

\section{A cleaner generating set}\label{sec.generators}

\Cref{con.head} bundles two very different things into one generator: three
\emph{linear} projections and one genuinely nonlinear routing step. It is
cleaner to separate them.

\begin{construction}[Linear part plus one routing primitive]\label{con.factor}
The projections $\Wq,\Wk,\Wv$ are linear arrangements, already living in
the linear suboperad of $\sarr$ (where the paper's symplectic-/linear
machinery applies). The one nonlinear primitive is row-softmax of the Gram
matrix,
\[
  \mathrm{route}\colon (\rr^{d})^{K}\longrightarrow (\Delta^{K-1})^{K},
  \qquad
  \mathrm{route}(u)_{i}=\softmax\bigl(\langle u_i,u_j\rangle\bigr)_{j},
\]
landing in a product of probability simplices. Then an attention head is
$\mathrm{linear}\circ\mathrm{route}\circ\mathrm{linear}$, and the only
object that needs separate characterization is $\mathrm{route}$.
\end{construction}

\begin{remark}[The linear fragment, categorified]\label{rmk.oneill}
O'Neill~\cite{ONeill} categorifies precisely this linear part: $\Wq,\Wk,\Wv$
as parametric morphisms in $\Cat{Para}(\vect)$ and layer-stacking as the free
monad on the induced endofunctor, with permutation equivariance and
mechanistic-interpretability circuits falling out as composites. His
$\Cat{Para}(\vect)$ is the forward leg of the parametric lens of gradient-based
learning~\cite{CGGWZ}; the backward leg those frameworks bolt on for training is
here intrinsic, the cotangent lift of $U$ pulled through the reactive sharp
$\sharpR_Q$ (\cref{rmk.params}). The softmax he
defers as needing ``more advanced categorical constructions'' is exactly the
primitive $\mathrm{route}$ isolated above; the operadic substitution
of~\cref{sec.suboperad} generalizes his single-type free-monad stacking to
arity-indexed trees, and the potential $U$ supplies dynamics for which a linear,
architecture-only reading has no room.
\end{remark}

\begin{remark}[Routing as data-dependent stochastic wiring]\label{rmk.stoch}
Row-softmax makes the attention matrix $A$ row-stochastic, so
$\mathrm{route}$ is a smooth section into stochastic matrices: a
data-dependent Markov kernel on the box set. The natural ambient operad is
that of \emph{data-dependent stochastic wiring diagrams}, with the
classical $0/1$ wiring diagrams as the deterministic corner. This is the
cleanest reading of ``the boxes learn a communication graph''; softmax
attention is one smooth generator of it.
\end{remark}

\begin{remark}[Parameters as state, learned by the same flow]\label{rmk.params}
Because the weights $W$ are the reactive space $\rv Q$ and the
configuration integrator's state space is $Q$ itself, $\Phiconf$ flows the
weights down $-\,\mathrm dU$ pulled through $\sharpR_Q$ --- the very same
gradient flow that relaxes the box values. Inference (relaxing $h$) and
learning (relaxing $W$) are one cotangent flow of one potential. Nothing
special needs to be said about ``training.''
\end{remark}

\section{Catalog of the generators we use}\label{sec.catalog}

For reference, the generators of the suboperad in one place, each presented as
the triple~\eqref{eq.arr}: the carrier $Q$ and its sharp $\sharpR_Q$, then the
three maps $\outp f,\inpt f,U$ of \eqref{eq.outf}--\eqref{eq.U}. Both legs of
every box are the residual stream $\rr^E$; $h\in\outp M$ is the inner-emission
tuple and $n\coloneqq\inpt N\in\rr^E$ the descending context. (The first two are
the paper's head and prior; the last two are route-2 extensions developed in the
companion notes --- a relaxable state carrier and a second-order read-out.)

\subsection*{Attention head (\cref{con.head}), arity $K\ge1$}
\[
\begin{aligned}
  Q          &= \vect\bigl(\rr^E,\rr^{d}\oplus\rr^{d}\oplus\rr^{d_v}\bigr)
                \oplus\vect\bigl(\rr^{d_v},\rr^E\bigr)\ni W=(\Wq,\Wk,\Wv,\Wo;w_\star),\\
  \sharpR_Q  &= \text{a metric on weight space (Euclidean by default, \cref{rmk.optimizer})},\\
  \outp f(W,h)     &= \Wo\sum_{j=1}^{K}\softmax\bigl(\tfrac1{\sqrt d}\langle w_\star,\Wk h_j\rangle\bigr)_j\,\Wv h_j,\\
  \inpt f(W,h,n)_i &= \Wo\!\!\sum_{j\in\{1,\dots,K\}\cup\{\star\}}\!\!A_{ij}\,\Wv h_j,
                     \qquad A_{i\cdot}=\softmax\bigl(\tfrac1{\sqrt d}\langle\Wq h_i,\Wk h_j\rangle\bigr),\\
  U(W,h,n)         &= \tfrac12\bigl\lVert\,h-\Wo\Wv n\,\bigr\rVert^{2}.
\end{aligned}
\]

\subsection*{Prior box (\cref{ex.zeroary}), arity $K=0$}
\[
\begin{aligned}
  Q              &= \rr^{E}\ni q,\\
  \sharpR_Q      &= \text{the Euclidean metric on }\rr^{E},\\
  \outp f(q,h)   &= q,\\
  \inpt f(q,h,n) &= \langle\rangle\quad(\text{the unique map to }\rr^{0}),\\
  U(q,h,n)       &= 0.
\end{aligned}
\]
Here $\outp M=\inpt M=\rr^{0}$, so $h$ is the empty tuple.

\subsection*{Activation cell (route 2), arity $K=1$}
\[
\begin{aligned}
  Q                    &= \rr^{k}\oplus\Hom\bigl(\rr^{k},\rr^E\bigr)\ni(z,D),\\
  \sharpR_Q            &= \text{a block sharp: }\eta_z\,\mathrm I\text{ on }z\text{ (fast)},\ \eta_D\,\mathrm I\text{ on }D\text{ (slow)},\\
  \outp f\bigl((z,D),h\bigr)     &= D\,z,\\
  \inpt f\bigl((z,D),h,n\bigr)   &= n,\\
  U\bigl((z,D),h,n\bigr)         &= \tfrac12\bigl\lVert\,h-D\,z\,\bigr\rVert^{2}.
\end{aligned}
\]

\subsection*{Coincidence head (route 2, second order), arity $K\ge1$}
\[
\begin{aligned}
  Q                    &= \vect\bigl(\rr^E,\rr^{r}\bigr)\oplus\vect\bigl(\operatorname{Sym}^2\rr^{r},\rr^E\bigr)
                          \oplus\vect\bigl(\rr^E,\operatorname{Sym}^2\rr^{r}\bigr)\ni(V,\Wo,G),\\
  \sharpR_Q            &= \text{a metric on }(V,\Wo,G)\text{ (Euclidean by default)},\\
  \outp f\bigl((V,\Wo,G),h\bigr)   &= \Wo\,\operatorname{vech}\bigl(P P^{\top}\bigr),\qquad P=\sum_{j=1}^{K}V h_j\in\rr^{r},\\
  \inpt f\bigl((V,\Wo,G),h,n\bigr)_i &= n,\\
  U\bigl((V,\Wo,G),h,n\bigr)       &= \tfrac12\bigl\lVert\,\operatorname{vech}\bigl(P P^{\top}\bigr)-G\,n\,\bigr\rVert^{2}.
\end{aligned}
\]
The pooled feature $P$ is first order; $P P^{\top}$ makes the cross terms
$(V h_i)(V h_j)^{\top}$ --- the coincidences between children --- explicit, so
$\outp f$ is a degree-two read-out and $U$ predicts the second moment from $n$.

\section{Caveats and open questions}\label{sec.open}

\begin{enumerate}
\item \emph{Predictive-coding-shaped, not literally Rao--Ballard.} The name
in \cref{slo.pc} is earned by the potential and the relaxation dynamics,
the parts a bare transformer drops. But the routing is attention --- richer
than the linear generative model of classical predictive coding. The flat
linear corner is classical PC, the flat softmax corner is Hopfield, and the
nested attention version is the object actually defined here.

\item \emph{Self-attention is a feedback wiring.} The diagonal $A_{ii}>0$
is a self-loop; an acyclic wiring operad would forbid it. It is well-posed
here only through the dynamical reading (relax to a fixed point), which is
again \cref{rmk.corners}.

\item \emph{What is $\outp N$?} We took the pooled, super-token read-out
of~\eqref{eq.pool} (a learned-query pool), which makes nesting hierarchical
(\cref{rmk.supertoken}). Passing the whole length-$K$ sequence upward
instead gives a different composability regime (segment/hierarchical
attention); the interface for substitution differs accordingly.

\item \emph{Characterize the generated suboperad.} \Cref{obs.generate}
asserts closure abstractly. Give an explicit normal form for finite
attention networks inside $\sarr$, and decide whether $\mathrm{route}$
(\cref{con.factor}) is the unique nonlinear generator needed --- i.e.\
whether the linear suboperad together with $\mathrm{route}$ generates
exactly the transformer-with-residual class once the output projection
$\Wo$ closes the residual stream (\cref{obs.residual}).

\item \emph{Is attention wiring or force?} In \cref{con.head} attention is
the routing $\inpt f$ and prediction is a separate potential $U$, so the
dynamics is a gradient flow of $U$ for any (untied) $\Wo$. In the energy
reading (\cref{rmk.corners}) attention \emph{is} $-\nabla U$, which forces
the symmetric, tied circuits of \cref{rmk.circuits}. Decide whether the
tied/symmetric head --- where routing and force coincide --- is the
canonical object and the untied transformer a deliberate weakening.

\item \emph{A Hopfield proposition.} State and prove that $\Phiconf$ of the
decorated head of \cref{con.head} is the energy-descent update of a modern
Hopfield network, making \cref{rmk.corners} a theorem rather than a
correspondence.

\item \emph{Close the loop.} Make \cref{rmk.active} precise: which
arrangement closes the down-leg onto the environment, and what
conservation/dissipation statement distinguishes the active-inference
dynamics from the perceptual one.
\end{enumerate}

\section{A talking operad, to be checked}\label{sec.talk}

Here is a provisional way to make precise the slogan that boxes can
learn to talk to one another. The point is to put a typed message
manifold between raw emissions and attention routing. Attention is then
communication factored through messages.

Fix a smooth language signature
\[
  \Lambda=(A,P),
\]
where $A$ is an address manifold and $P$ is a payload manifold. Define
the message manifold
\[
  \mathrm{Msg}_{\Lambda}
    \coloneqq A_{\mathrm{to}}\times P\times A_{\mathrm{return}} .
\]
A point of $\mathrm{Msg}_{\Lambda}$ says: send payload $p$ to address
$a_{\mathrm{to}}$, and replies should go to address
$a_{\mathrm{return}}$.

Define a colored operad $\mathsf{Talk}_{\Lambda}$ as follows. Its colors
are the same lens objects $\lensob X=\binom{\inpt X}{\outp X}$ as in
$\sarr$. A one-round talking generator
\[
  \tau\colon \lensob{X_1},\ldots,\lensob{X_K}\longrightarrow\lensob Y
\]
consists of the following smooth data.

\begin{enumerate}
\item A reactive parameter space $\rv Q=(Q,\sharpR_Q)$.

\item Speaker maps, one for each internal box,
\[
  \mu_j\colon Q\times\outp{X_j}\longrightarrow \mathrm{Msg}_{\Lambda}
  \qquad(1\leq j\leq K),
\]
and a descending speaker map for the outer input,
\[
  \mu_\star\colon Q\times\inpt Y\longrightarrow\mathrm{Msg}_{\Lambda}.
\]

\item Listener-address maps
\[
  \lambda_i\colon Q\times\outp{X_i}\longrightarrow A
  \qquad(1\leq i\leq K),
\]
and an outer listener-address map
\[
  \lambda_\star\colon Q\times\outp{X_1}\times\cdots\times\outp{X_K}
    \longrightarrow A .
\]
The latter has no $\inpt Y$-argument, preserving the Moore condition.

\item A learned compatibility score
\[
  \kappa\colon Q\times A\times A\longrightarrow\rr .
\]

\item Decoder maps
\[
  D_i\colon
  Q\times\outp{X_i}\times\Delta^K_{\circ}
  \times\mathrm{Msg}_{\Lambda}^{K+1}
  \longrightarrow\inpt{X_i}
  \qquad(1\leq i\leq K),
\]
and
\[
  D_\star\colon
  Q\times\Delta^{K-1}_{\circ}\times\mathrm{Msg}_{\Lambda}^{K}
  \longrightarrow\outp Y .
\]
Here $\Delta^K_{\circ}$ denotes the open $K$-simplex, so the routing
weights remain smooth manifold-valued.

\item A language potential
\[
  V\colon Q\times\mathrm{Msg}_{\Lambda}^{K}\times\mathrm{Msg}_{\Lambda}
    \longrightarrow\rr .
\]
\end{enumerate}

Given internal emissions $h=(h_1,\ldots,h_K)$ and an outer input
$n\in\inpt Y$, put
\[
  m_j\coloneqq\mu_j(q,h_j),
  \qquad
  m_\star\coloneqq\mu_\star(q,n),
\]
and write $a_j^{\mathrm{to}}$ and $a_\star^{\mathrm{to}}$ for their
destination-address components. The internal routing weights are
\[
  \alpha_{ij}
  =
  \softmax_{j\in\{1,\ldots,K,\star\}}
  \kappa\bigl(q,\lambda_i(q,h_i),a_j^{\mathrm{to}}\bigr),
\]
and the outer routing weights are
\[
  \alpha_{\star j}
  =
  \softmax_{j\in\{1,\ldots,K\}}
  \kappa\bigl(q,\lambda_\star(q,h),a_j^{\mathrm{to}}\bigr).
\]

There is then an evident assignment on generators
\[
  F_\Lambda\colon\mathsf{Talk}_{\Lambda}\longrightarrow\sarr
\]
sending $\tau$ to the smooth adaptive arrangement
\[
  F_\Lambda(\tau)=(\rv Q,\outp f,\inpt f,U)
\]
defined by
\[
  \outp f(q,h)
  =
  D_\star\bigl(q,(\alpha_{\star j})_j,(m_j)_j\bigr),
\]
\[
  \inpt f(q,h,n)_i
  =
  D_i\bigl(q,h_i,(\alpha_{ij})_{j\in\{1,\ldots,K,\star\}},
    (m_1,\ldots,m_K,m_\star)\bigr),
\]
and
\[
  U(q,h,n)=V(q,m_1,\ldots,m_K,m_\star).
\]
The operad $\mathsf{Talk}_{\Lambda}$ should be defined as the colored
operad freely generated by these one-round talking arrangements, with the
usual symmetric actions on internal boxes. The assignment above then
extends uniquely to an operad functor
\[
  F_\Lambda\colon\mathsf{Talk}_{\Lambda}\to\sarr,
\]
using operadic substitution in $\sarr$.

\begin{slogan}
A learned language is attention routing factored through a message
manifold:
\[
  \outp X
  \longrightarrow
  \mathrm{Msg}_{\Lambda}
  \longrightarrow
  \Delta\times\text{payloads}.
\]
\end{slogan}

In this reading, the return-address component is not extra syntax beyond
the message type. A box answers a question by later emitting a message
whose destination address is the return address it received. The maps
$\mu,\lambda,\kappa,D$ and the potential $V$ all live in the reactive
parameter space, so under $\Phiconf$ the language itself is learned by the
same flow as the rest of the arrangement.

%=====================================================================
\begin{thebibliography}{9}

\bibitem{DAP} (Companion manuscript.) \emph{Dynamic algebra of potentials},
\texttt{dynamic-algebra-potentials.tex}. Definitions of the smooth
arrangement operad $\sarr$, the cotangent interpretation $\cot$, and the
integrators $\Phiconf,\Phiphase$ are taken from there.

\bibitem{HopfieldAllYou} H.~Ramsauer et al.,
\emph{Hopfield Networks is All You Need}, ICLR 2021.

\bibitem{RaoBallard} R.~P.~N.~Rao and D.~H.~Ballard,
\emph{Predictive coding in the visual cortex}, Nature Neuroscience 2(1), 1999.

\bibitem{Friston} K.~Friston,
\emph{The free-energy principle: a unified brain theory?},
Nature Reviews Neuroscience 11, 2010.

\bibitem{Attention} A.~Vaswani et al.,
\emph{Attention Is All You Need}, NeurIPS 2017.

\bibitem{SetTransformer} J.~Lee et al.,
\emph{Set Transformer}, ICML 2019.

\bibitem{Perceiver} A.~Jaegle et al.,
\emph{Perceiver: General Perception with Iterative Attention}, ICML 2021.

\bibitem{ONeill} C.~O'Neill,
\emph{Self-Attention as a Parametric Endofunctor: A Categorical Framework for
Transformer Architectures}, arXiv:2501.02931, 2025.

\bibitem{CGGWZ} G.~S.~H.~Cruttwell, B.~Gavranovi\'{c}, N.~Ghani, P.~Wilson,
and F.~Zanasi, \emph{Categorical Foundations of Gradient-Based Learning},
ESOP 2022, arXiv:2103.01931.

\end{thebibliography}

\end{document}
